\sezione{7}{Generalizzazioni}{slide 125–144}

\begin{inbreve}
Una \hl{generalizzazione} lega un'entità «genitore» a una o più entità «figlie» che ne sono \key{casi particolari}: ogni professore ordinario è anche un docente. Le figlie \key{ereditano} tutte le proprietà del genitore. Le generalizzazioni sono \hl{totali} o \hl{parziali}, \hl{esclusive} o \hl{sovrapposte}; nel corso si usano solo quelle esclusive.
\end{inbreve}

\subsection{Definizione e notazione}

\begin{tecnica}[Generalizzazione]
Mette in relazione una o più entità $E_1, E_2, \dots, E_n$ con un'entità $E$ che le comprende come casi particolari:
\begin{itemize}
  \item $E$ è \hl{generalizzazione} di $E_1, E_2, \dots, E_n$ (entità \key{genitore});
  \item $E_1, E_2, \dots, E_n$ sono \hl{specializzazioni} (o sottotipi) di $E$ (entità \key{figlie}).
\end{itemize}
\end{tecnica}

\par\noindent\begin{minipage}{\linewidth}
\begin{schema}[x=1cm,y=1cm]
  \node[ent] (D) at (0,0) {Docente};
  \node[ent] (O) at (-3.4,-2.2) {Professore\\ordinario};
  \node[ent] (A) at (0,-2.2) {Professore\\associato};
  \node[ent] (R) at (3.4,-2.2) {Ricercatore};
  \draw[line width=.9pt,navy] (O.north) |- (0,-1.15) (R.north) |- (0,-1.15) (A.north) -- (0,-1.15);
  \draw[genp] (0,-1.15) -- (D.south);
\end{schema}
\didascalia{Rappresentazione grafica di una generalizzazione (slide 128): le figlie sono collegate al genitore da una freccia.}
\end{minipage}\par

\subsection{Caratteristiche ed ereditarietà}

\begin{tecnica}[Proprietà delle generalizzazioni]
Se $E$ (genitore) è generalizzazione di $E_1, \dots, E_n$ (figlie):
\begin{itemize}
  \item ogni occorrenza di $E_1, \dots, E_n$ è \hl{anche occorrenza di $E$};
  \item ogni proprietà di $E$ è significativa per $E_1, \dots, E_n$.
\end{itemize}
\hl{Ereditarietà}: tutte le proprietà dell'entità genitore (attributi, relationship, altre generalizzazioni, identificatori) vengono ereditate dalle figlie e \key{non si rappresentano esplicitamente}.
\end{tecnica}

\par\noindent\begin{minipage}{\linewidth}
\begin{schema}[x=1cm,y=1cm]
  \node[ent] (P) at (0,0) {Persona};
  \node[rel] (N) at (0,1.75) {Nascita};
  \node[ent] (C) at (0,3.5) {Città};
  \draw[lk] (P) -- (N) node[card,pos=.3,right]{(1,1)};
  \draw[lk] (N) -- (C) node[card,pos=.75,right]{(0,N)};
  \idattributo{(P.east)++(0,.2)}{(2.6,.65)}{Codice fiscale}{right}
  \attributo{(P.east)}{(2.6,0)}{Nome}{right}
  \attributo{(P.east)++(0,-.2)}{(2.6,-.65)}{Età}{right}
  \node[ent] (L) at (-2.4,-2.3) {Lavoratore};
  \node[ent] (S) at (2.4,-2.3) {Studente};
  \draw[line width=.9pt,navy] (L.north) |- (0,-1.3) -| (S.north);
  \draw[genp] (0,-1.3) -- (P.south);
  \attributo{(L.west)}{(-4.6,-2.3)}{Stipendio}{left}
\end{schema}
\didascalia{Generalizzazione ed ereditarietà (slide 130): Lavoratore e Studente hanno anche codice fiscale, nome, età e città di nascita; Lavoratore ha in più lo stipendio.}
\end{minipage}\par

\begin{esempio}[Che cosa eredita uno studente]
Uno studente è una persona: ha quindi \hl{codice fiscale, nome ed età} ed è nato in \hl{una città}. Lo schema non ripete questi attributi su Studente: li eredita. Anche l'identificatore si eredita: uno studente è identificato dal codice fiscale.
\end{esempio}

\subsection{Tipi di generalizzazione}

\begin{tecnica}[Totale/parziale, esclusiva/sovrapposta]
\begin{itemize}
  \item \hl{Totale} se ogni occorrenza del genitore è occorrenza di \key{almeno una} figlia; altrimenti \hl{parziale}.
  \item \hl{Esclusiva} se ogni occorrenza del genitore è occorrenza di \key{al più una} figlia; altrimenti \hl{sovrapposta}.
\end{itemize}
Nel corso si considerano, senza perdita di generalità, solo generalizzazioni \key{esclusive}, distinguendo fra totali e parziali. Graficamente: \key{freccia vuota} = parziale, \key{freccia piena} = totale.
\end{tecnica}

\par\noindent\begin{minipage}{\linewidth}
\begin{schema}[x=1cm,y=1cm]
  \node[ent] (P) at (0,0) {Persona};
  \node[ent] (Di) at (-1.9,-2.2) {Disoccupato};
  \node[ent] (La) at (1.9,-2.2) {Lavoratore};
  \draw[line width=.9pt,navy] (Di.north) |- (0,-1.2) -| (La.north);
  \draw[genp] (0,-1.2) -- (P.south);
  \node[nota] at (0,-3.15) {freccia vuota: parziale};
  \node[ent] (Pa) at (8,0) {Parlamentare};
  \node[ent] (De) at (6.1,-2.2) {Deputato};
  \node[ent] (Se) at (9.9,-2.2) {Senatore};
  \draw[line width=.9pt,navy] (De.north) |- (8,-1.2) -| (Se.north);
  \draw[gent] (8,-1.2) -- (Pa.south);
  \node[nota] at (8,-3.15) {freccia piena: totale};
\end{schema}
\didascalia{Generalizzazione parziale (ci sono persone né disoccupate né lavoratrici, per esempio gli studenti) e totale (ogni parlamentare è deputato o senatore), slide 133–134.}
\end{minipage}\par

\begin{esempio}[Leggere le due proprietà con gli insiemi]
Indicando con $E$ l'insieme delle occorrenze del genitore e con $E_i$ quelli delle figlie:
\[
\text{totale:}\ \ E = E_1\cup\dots\cup E_n \qquad\qquad \text{esclusiva:}\ \ E_i\cap E_j=\varnothing \ \ \text{per ogni}\ i\neq j .
\]
Parlamentare, Deputato, Senatore: ogni parlamentare appartiene a una delle due camere (totale) e a una sola (esclusiva).
\end{esempio}

\subsection{Altre proprietà}

\begin{tecnica}[Gerarchie, sottoinsiemi, identificatori]
\begin{itemize}
  \item Possono esistere \hl{gerarchie a più livelli} e più generalizzazioni allo stesso livello.
  \item Un'entità può far parte di più gerarchie, come genitore e/o come figlia.
  \item Una generalizzazione con \hl{una sola figlia} si chiama \key{sottoinsieme}.
  \item Il genitore di una generalizzazione \hl{totale} può non avere identificatore, purché ogni figlia ne abbia uno (risposta ricostruita alla domanda della slide 138): ogni occorrenza del genitore è occorrenza di una figlia e lì è identificata.
\end{itemize}
\end{tecnica}

\par\noindent\begin{minipage}{\linewidth}
\begin{schema}[x=1cm,y=1cm]
  \node[ent,minimum width=1.4cm] (E0) at (0,0) {E0};
  \node[ent,minimum width=1.4cm] (E1) at (-3.6,-1.8) {E1};
  \node[ent,minimum width=1.4cm] (E2) at (-1.2,-1.8) {E2};
  \node[ent,minimum width=1.4cm] (E6) at (1.8,-1.8) {E6};
  \node[ent,minimum width=1.4cm] (E7) at (4.0,-1.8) {E7};
  \node[ent,minimum width=1.4cm] (E3) at (-4.6,-3.6) {E3};
  \node[ent,minimum width=1.4cm] (E4) at (-2.6,-3.6) {E4};
  \node[ent,minimum width=1.4cm] (E8) at (-1.2,-3.6) {E8};
  \draw[line width=.9pt,navy] (E1.north) |- (-2.4,-.95) -| (E2.north);
  \draw[genp] (-2.4,-.95) -- (-2.4,-.6) -- (-.3,-.6) -- (-.3,-.41);
  \draw[line width=.9pt,navy] (E6.north) |- (2.9,-.95) -| (E7.north);
  \draw[genp] (2.9,-.95) -- (2.9,-.6) -- (.3,-.6) -- (.3,-.41);
  \draw[line width=.9pt,navy] (E3.north) |- (-3.6,-2.75) -| (E4.north);
  \draw[gent] (-3.6,-2.75) -- (E1.south);
  \draw[genp] (E8.north) -- (E2.south);
\end{schema}
\didascalia{Gerarchie a più livelli e generalizzazioni multiple (slide 137): E0 ha due generalizzazioni distinte; E8 è un sottoinsieme di E2.}
\end{minipage}\par

\subsection{Configurazioni che non hanno senso}

\begin{attenzione}[Studenti e lavoratori]
\emph{Persona} con figlie \emph{Studente} e \emph{Lavoratore} in una generalizzazione \hl{esclusiva} non va bene se esistono studenti che lavorano. Non ha senso nemmeno mettere Studente come figlia di Lavoratore (o viceversa): dipende dal punto di vista («un'azienda con anche lavoratori studenti» oppure «un'università con anche studenti lavoratori») e in entrambi i casi esclude qualcuno. E un'entità \hl{figlia di due generalizzazioni} diverse (E3 figlia sia di E1 sia di E2) non è ammessa.
\end{attenzione}

\par\noindent\begin{minipage}{\linewidth}
\begin{schema}[x=1cm,y=1cm]
  \node[ent] (P) at (0,0) {Persona};
  \node[ent] (S) at (-2.6,-1.7) {Studente};
  \node[ent] (L) at (2.6,-1.7) {Lavoratore};
  \node[ent] (SL) at (0,-3.4) {Studente\\lavoratore};
  \draw[genp] (S.north) -- (P.south west);
  \draw[genp] (L.north) -- (P.south east);
  \draw[genp] (SL.west) -| (S.south);
  \draw[genp] (SL.east) -| (L.south);
\end{schema}
\didascalia{Una soluzione (slide 142): Studente e Lavoratore sono due sottoinsiemi distinti di Persona; Studente lavoratore è sottoinsieme di entrambi. Ogni freccia è una generalizzazione con una sola figlia, quindi non ci sono generalizzazioni sovrapposte.}
\end{minipage}\par

\subsection{Esercizio guidato}

\begin{esempio}[Persone, impiegati e studenti (slide 143–144)]
\emph{Le persone hanno CF, cognome ed età; gli uomini anche la posizione militare; gli impiegati hanno lo stipendio e possono essere segretari, direttori o progettisti (un progettista può essere anche responsabile di progetto); gli studenti (che non possono essere impiegati) un numero di matricola; esistono persone che non sono né impiegati né studenti.}
\begin{itemize}
  \item Persona → \{Uomo, Donna\}: \hl{totale} ed esclusiva; Uomo ha l'attributo Militare.
  \item Persona → \{Impiegato, Studente\}: \hl{parziale} (ci sono persone né impiegate né studenti) ed esclusiva (gli studenti non possono essere impiegati).
  \item Impiegato → \{Segretario, Direttore, Progettista\}: \hl{totale} («possono essere» tre categorie; nello schema delle slide è disegnata piena).
  \item Progettista → Responsabile: \hl{sottoinsieme}.
\end{itemize}
\end{esempio}

\par\noindent\begin{minipage}{\linewidth}
\begin{schema}[x=1cm,y=1cm,scale=.92,every node/.append style={transform shape}]
  \node[ent,minimum width=5.6cm] (P) at (0,0) {Persona};
  \idattributo{(P.north west)++(.3,0)}{(-2.5,1.1)}{CF}{left}
  \attributo{(P.west)}{(-4.2,0)}{Cognome}{left}
  \attributo{(P.east)}{(4.2,0)}{Età}{right}
  \node[ent,minimum width=1.6cm] (U) at (-5.2,-2) {Uomo};
  \node[ent,minimum width=1.6cm] (D) at (-3.2,-2) {Donna};
  \node[ent] (I) at (0.4,-2) {Impiegato};
  \node[ent] (S) at (3.8,-2) {Studente};
  \draw[line width=.9pt,navy] (U.north) |- (-4.2,-1.15) -| (D.north);
  \draw[gent] (-4.2,-1.15) -- (-4.2,-.8) -- (-1.6,-.8) -- (-1.6,-.41);
  \draw[line width=.9pt,navy] (I.north) |- (2.1,-1.15) -| (S.north);
  \draw[genp] (2.1,-1.15) -- (2.1,-.41);
  \attributo{(U.south)}{(-5.2,-3.1)}{Militare}{below}
  \attributo{(I.east)++(0,-.2)}{(2.0,-2.85)}{Stipendio}{right}
  \attributo{(S.east)}{(5.8,-2)}{Matricola}{right}
  \node[ent] (Sg) at (-2.6,-4.1) {Segretario};
  \node[ent] (Dr) at (0.4,-4.1) {Direttore};
  \node[ent] (Pr) at (3.4,-4.1) {Progettista};
  \draw[line width=.9pt,navy] (Sg.north) |- (.4,-3.2) -| (Pr.north) (Dr.north) -- (.4,-3.2);
  \draw[gent] (.4,-3.2) -- (I.south);
  \node[ent] (Re) at (3.4,-5.9) {Responsabile};
  \draw[genp] (Re.north) -- (Pr.south);
\end{schema}
\didascalia{Soluzione dell'esercizio (slide 144).}
\end{minipage}\par

\begin{esercizi}
\begin{passi}
  \item Gli utenti di una palestra sono soci o ospiti giornalieri; i soci possono essere anche istruttori. Disegna la gerarchia indicando totalità ed esclusività.
  \item In una generalizzazione totale il genitore \emph{Veicolo} non ha identificatore. Quale condizione deve valere sulle figlie \emph{Auto} e \emph{Moto}?
\end{passi}
\end{esercizi}

\begin{appunti}{Riepilogo del capitolo}
\pag{127} Genitore $E$ e figlie $E_1, \dots, E_n$ (specializzazioni).
\pag{129} Ogni occorrenza di una figlia è occorrenza del genitore.
\pag{131} Ereditarietà: attributi, relationship, identificatori del genitore valgono per le figlie.
\pag{132} Totale/parziale (freccia piena/vuota); esclusiva/sovrapposta (nel corso solo esclusive).
\pag{136} Gerarchie a più livelli, generalizzazioni multiple, sottoinsiemi.
\pag{142} Studenti lavoratori: due sottoinsiemi e un sottoinsieme comune.
\end{appunti}
